\section{4.6.SWC Catalog}\label{swc-catalog}

\subsection{4.6. Catálogo de vulnerabilidades
SWC}\label{catuxe1logo-de-vulnerabilidades-swc}

Con el objetivo de desacoplar la fase de detección de vulnerabilidades
de las fases posteriores de validación y generación de pruebas, se
implementó un catálogo de vulnerabilidades basado en la clasificación
Smart Contract Weakness Classification (SWC).

Este catálogo constituye una capa intermedia que permite asociar los
hallazgos detectados con información adicional independiente de las
herramientas utilizadas.

Gracias a ello, EVMAudit puede reutilizar la información obtenida
durante el proceso de correlación y emplearla posteriormente para la
generación automática de propiedades orientadas a Echidna.

\subsubsection{4.6.1. Objetivos del
catálogo}\label{objetivos-del-catuxe1logo}

El catálogo desarrollado persigue los siguientes objetivos:

\begin{itemize}
\tightlist
\item
  unificar la representación de vulnerabilidades
\item
  desacoplar la lógica de generación de pruebas
\item
  facilitar la incorporación de nuevas plantillas
\item
  centralizar información sobre las debilidades conocidas
\item
  favorecer la mantenibilidad del sistema.
\end{itemize}

Esta aproximación permite que las fases posteriores del análisis no
dependan directamente de las particularidades de Slither o Mythril.

\subsubsection{4.6.2. Estructura del
catálogo}\label{estructura-del-catuxe1logo}

Cada entrada del catálogo contiene información asociada a una
vulnerabilidad concreta.

Entre los datos almacenados se encuentran:

\begin{itemize}
\tightlist
\item
  identificador SWC
\item
  nombre de la vulnerabilidad
\item
  descripción
\item
  severidad
\item
  plantilla asociada
\item
  limitaciones conocidas
\item
  posibilidad de validación automática.
\end{itemize}

La información almacenada permite enriquecer los hallazgos
correlacionados y proporcionar información adicional durante las fases
posteriores del análisis.

!TODO: insertar tabla con varios ejemplos del catálogo.

\subsubsection{4.6.3. Asociación entre detectores y
SWC}\label{asociaciuxf3n-entre-detectores-y-swc}

Una de las principales dificultades observadas durante el desarrollo fue
la ausencia de una correspondencia directa entre los detectores
utilizados por cada herramienta.

Para resolver este problema se definieron tablas de asociación que
permiten traducir los detectores específicos a identificadores SWC.

Gracias a esta estrategia es posible utilizar un lenguaje común durante
todo el proceso de análisis y evitar dependencias directas respecto a la
nomenclatura propia de cada herramienta.

Esta decisión resulta especialmente relevante para el módulo de
correlación descrito anteriormente.

!TODO: insertar referencia cruzada a la Sección 4.4.

\subsubsection{4.6.4. Reutilización del
catálogo}\label{reutilizaciuxf3n-del-catuxe1logo}

Además de servir como mecanismo de clasificación, el catálogo es
utilizado posteriormente durante la generación de propiedades para
Echidna.

A partir del identificador SWC asociado a cada vulnerabilidad, EVMAudit
recupera automáticamente la plantilla correspondiente y genera las
estructuras necesarias para la fase de fuzzing.

Esta aproximación permite separar claramente:

\begin{itemize}
\tightlist
\item
  detección
\item
  correlación
\item
  generación de pruebas.
\end{itemize}

La independencia entre estas fases facilita futuras ampliaciones y
simplifica el mantenimiento del sistema.

\subsubsection{4.6.5. Extensibilidad del
catálogo}\label{extensibilidad-del-catuxe1logo}

La utilización de un catálogo independiente permite incorporar nuevas
vulnerabilidades sin necesidad de modificar el resto de módulos
implementados.

De esta forma, la incorporación de nuevas entradas únicamente requiere:

\begin{enumerate}
\def\labelenumi{\arabic{enumi}.}
\tightlist
\item
  definir el identificador SWC correspondiente
\item
  añadir la información descriptiva asociada
\item
  incorporar la plantilla de generación deseada.
\end{enumerate}

Este diseño favorece la evolución futura de la herramienta y permite
adaptar EVMAudit a nuevas versiones del ecosistema Ethereum.

!TODO: insertar diagrama de relación entre SWC Catalog y Echidna
Adapter.

!TODO: insertar referencia a la Figura X.
